%% Lecture 13 -- Mean Value Theorem; Dirichlet & Neumann Green's Functions; Method of Images
\section{Lecture 13 -- Green's Functions for Boundary-Value Problems}\label{sec:lec13}

\begin{center}
\begin{tabular}{ll}
  \textbf{Date:}\quad 19/05/2026 & \\
\end{tabular}
\end{center}
\hrule\vspace{1.5em}

\subsection*{Overview}

This lecture completes the programme begun in Lecture 12.  We have a representation
formula for $\phi$ in terms of boundary data, but it is \emph{not} a solution because it
contains both $\phi|_\Sigma$ and $\partial_n\phi|_\Sigma$ simultaneously.  We:
\begin{itemize}
  \item prove the \textbf{mean value property} of harmonic functions from the representation
        formula, and use it to establish Earnshaw's theorem rigorously;
  \item generalize from the bare $1/|\br-\mathbf{y}|$ to an arbitrary Green's function $G$,
        deriving the general representation formula;
  \item define the \textbf{Dirichlet} and \textbf{Neumann} Green's functions and show
        how each eliminates the unwanted boundary term;
  \item expose the inconsistency of the naive Neumann condition $\partial_n G_N = 0$ and
        derive the correct condition;
  \item connect the whole framework to the \textbf{method of images} studied in earlier courses;
  \item state the symmetry property $G_D(\br,\mathbf{y}) = G_D(\mathbf{y},\br)$
        (proof deferred to next lecture).
\end{itemize}

%------------------------------------------------------------
\subsection{Proof of the Mean Value Property}\label{subsec:lec13-mvp}

\subsubsection*{Motivation}
Last lecture we proved (but did not use) Earnshaw's theorem by appealing to the maximum
principle.  Here we derive it rigorously from the representation formula.  The key tool is
the \emph{mean value property}: for a harmonic function $u$ (one satisfying
$\nabla^2 u = 0$), the value at any interior point equals the average over any sphere
centred there.

\subsubsection*{Setup}
Let $u$ satisfy Laplace's equation in a domain $\Omega$:
\begin{equation}\label{eq:lec13:laplace}
  \nabla^2 u = 0 \quad \text{in } \Omega.
\end{equation}
Pick any interior point $\mathbf{y}\in\Omega$ and any sphere $B_R(\mathbf{y})$ of radius $R$
that lies entirely inside $\Omega$.  Apply the representation formula (derived last lecture)
\emph{to this smaller domain} $B_R(\mathbf{y})$, with $\psi(\br;\mathbf{y}) = 1/|\br-\mathbf{y}|$:

\begin{equation}\label{eq:lec13:rep-harmonic}
  u(\mathbf{y})
  = -\frac{1}{4\pi}
  \oint_{\partial B_R}
  \!\!\left[
    u(\br)\,\partial_n\!\left(\frac{1}{|\br-\mathbf{y}|}\right)
    - \frac{1}{|\br-\mathbf{y}|}\,\partial_n u(\br)
  \right] dA.
\end{equation}

On the sphere $|\br - \mathbf{y}| = R$, so $\psi = 1/R$ is constant.
The normal $\hat{n}$ points radially outward, so
\begin{equation}\label{eq:lec13:psi-deriv}
  \partial_n\!\left(\frac{1}{|\br-\mathbf{y}|}\right)\bigg|_{|\br-\mathbf{y}|=R}
  = -\frac{1}{R^2}.
\end{equation}

Substituting into \eqref{eq:lec13:rep-harmonic}:
\begin{equation}\label{eq:lec13:two-terms}
  u(\mathbf{y})
  = -\frac{1}{4\pi}
  \oint_{\partial B_R}
  \!\!\left[
    u(\br)\cdot\!\left(-\frac{1}{R^2}\right)
    - \frac{1}{R}\,\partial_n u(\br)
  \right] dA
  = \frac{1}{4\pi R^2}\oint_{\partial B_R} u\,dA
    + \frac{1}{4\pi R}\oint_{\partial B_R} \partial_n u\,dA.
\end{equation}

The second term is $\frac{1}{4\pi R}\oint \nabla u \cdot \hat{n}\, dA$.  By Gauss's theorem
this equals $\frac{1}{4\pi R}\int_{B_R}\nabla^2 u\,d^3r = 0$ (since $u$ is harmonic).
Hence:

\begin{equation}\label{eq:lec13:mvp}
  \boxed{
    u(\mathbf{y}) = \frac{1}{4\pi R^2}\oint_{|\br-\mathbf{y}|=R} u(\br)\,dA
    = \langle u \rangle_{\partial B_R(\mathbf{y})}.
  }
\end{equation}
\textit{Significance:} The value of a harmonic function at any point equals its spherical
average over any concentric sphere inside the domain.  This is a global constraint
that a harmonic function must satisfy at \emph{every} interior point.

\begin{remark}
  The mean value \eqref{eq:lec13:mvp} holds for \emph{every} radius $R$ such that
  $B_R(\mathbf{y})\subset\Omega$.  It is therefore a powerful intrinsic property of
  harmonic functions, independent of the boundary geometry.
\end{remark}

%------------------------------------------------------------
\subsection{Earnshaw's Theorem from the Mean Value Property}\label{subsec:lec13-earnshaw}

\subsubsection*{Motivation}
A local minimum of $u$ at $\mathbf{y}$ means $u(\mathbf{y}) < u(\br)$ for all $\br$ in some
punctured neighbourhood.  But then the spherical average over a small enough sphere is
\emph{strictly greater} than $u(\mathbf{y})$, contradicting \eqref{eq:lec13:mvp}.

\thm[thm:earnshaw]{Earnshaw / Maximum Principle}{
  A harmonic function on a bounded domain $\Omega$ (satisfying $\nabla^2 u=0$)
  cannot attain a local interior minimum or maximum unless $u$ is constant.
  Extrema occur only on the boundary $\partial\Omega$.
}

\textit{Physical consequence:} No static arrangement of point charges can be in stable
equilibrium under purely electrostatic forces.  The potential $\Phi$ created by all
other charges satisfies $\nabla^2\Phi=0$ near any charge-free point, so $\Phi$ has no
interior minimum there — stable equilibrium is impossible.

%------------------------------------------------------------
\subsection{Generalising the Green's Function}\label{subsec:lec13-general-G}

\subsubsection*{Conceptual Background}
In the representation formula \eqref{eq:lec13:rep-harmonic} we used
$\psi = 1/|\br-\mathbf{y}|$, which satisfies
$\nabla^2_\br\psi = -4\pi\delta^{(3)}(\br-\mathbf{y})$.
\emph{Only this property was used.}  Any function with the same singularity structure
yields an equally valid formula.  This gives us a free function to play with — and we
will use it to eliminate one of the two boundary terms.

\dfn[dfn:green]{Green's Function of the Laplacian}{
  A function $G(\br,\mathbf{y})$ is called a \emph{Green's function for the Laplacian}
  if it satisfies
  \begin{equation}\label{eq:lec13:G-def}
    \nabla^2_\br\, G(\br,\mathbf{y}) = -4\pi\,\delta^{(3)}(\br - \mathbf{y}).
  \end{equation}
  The most general such function is
  \begin{equation}\label{eq:lec13:G-general}
    G(\br,\mathbf{y}) = \frac{1}{|\br-\mathbf{y}|} + F(\br,\mathbf{y}),
  \end{equation}
  where $F(\br,\mathbf{y})$ is \emph{harmonic in $\br$}: $\nabla^2_\br F = 0$ inside $\Omega$.
}

\subsubsection*{General Representation Formula}
Substituting $G$ (instead of $\psi = 1/|\br-\mathbf{y}|$) into Green's second identity and
using $\nabla^2\phi = -4\pi\rho$ (Poisson's equation) gives:

\begin{equation}\label{eq:lec13:rep-general}
  \boxed{
    \phi(\mathbf{y})
    = \int_\Omega \rho(\br)\, G(\br,\mathbf{y})\, d^3r
    \;-\; \frac{1}{4\pi}\oint_\Sigma
    \!\!\left[\phi(\br)\,\partial_n G(\br,\mathbf{y})
      - G(\br,\mathbf{y})\,\partial_n\phi(\br)\right]dA.
  }
\end{equation}
\textit{Significance:} This formula is valid for \emph{any} choice of $G$ satisfying
\eqref{eq:lec13:G-def}. The freedom in $F$ lets us impose one additional boundary
condition on $G$ so that one of the surface terms vanishes.

\begin{figure}[h]
\centering
\begin{tikzpicture}[>=Stealth, thick]
  % Domain blob
  \draw[blue!60!black, fill=blue!5, thick]
    plot[smooth cycle, tension=0.7]
    coordinates {(0,2.2)(1.6,2.8)(3.2,2.0)(3.6,0.8)(2.8,-0.4)(1.0,-0.6)(-0.4,0.4)};
  % Label domain
  \node at (1.7,1.0) {$\Omega$};
  % Normal vector on boundary
  \draw[->, red!70!black] (3.3,1.4) -- (4.0,1.8) node[right] {$\hat{n}$};
  % Point y inside
  \fill (1.6,1.6) circle (2.5pt) node[above right] {$\mathbf{y}$};
  % Small sphere around y
  \draw[dashed, gray] (1.6,1.6) circle (0.6);
  \node[gray] at (1.6, 2.35) {$B_R(\mathbf{y})$};
  % Boundary label
  \node at (0.5,2.2) {$\Sigma = \partial\Omega$};
  % phi=0 label for Dirichlet
  \node[blue!60!black] at (3.6,2.3) {$G_D|_\Sigma = 0$};
  % point charge symbol
  \node[below left] at (1.6,1.6) {\small $q{=}1$};
\end{tikzpicture}
\caption{Domain $\Omega$ with outward normal $\hat{n}$ on boundary $\Sigma$.  The
Dirichlet Green's function $G_D(\br,\mathbf{y})$ is the potential of a unit point charge
at $\mathbf{y}$ with the boundary held at zero — the method-of-images setup.
The dashed sphere $B_R(\mathbf{y})$ is used in the mean value proof.}
\label{fig:lec13:domain}
\end{figure}

%------------------------------------------------------------
\subsection{The Dirichlet Green's Function}\label{subsec:lec13-dirichlet}

\subsubsection*{Motivation}
In the Dirichlet problem we \emph{know} $\phi|_\Sigma = f(\br)$ but \emph{do not know}
$\partial_n\phi|_\Sigma$.  We want the term $G\,\partial_n\phi$ in \eqref{eq:lec13:rep-general}
to vanish.  The simplest way: require $G$ itself to vanish on $\Sigma$.

\dfn[dfn:green-dirichlet]{Dirichlet Green's Function $G_D$}{
  $G_D(\br,\mathbf{y})$ is the unique solution of
  \begin{equation}\label{eq:lec13:GD-def}
    \nabla^2_\br\, G_D(\br,\mathbf{y}) = -4\pi\,\delta^{(3)}(\br - \mathbf{y}),
    \qquad
    G_D(\br,\mathbf{y})\big|_{\br\in\Sigma} = 0.
  \end{equation}
}

With $G = G_D$, the surface term $G_D\,\partial_n\phi$ vanishes identically, and
\eqref{eq:lec13:rep-general} becomes the \textbf{Dirichlet solution formula}:

\begin{equation}\label{eq:lec13:dirichlet-solution}
  \boxed{
    \phi(\mathbf{y})
    = \int_\Omega \rho(\br)\, G_D(\br,\mathbf{y})\, d^3r
    \;-\; \frac{1}{4\pi}\oint_\Sigma
    f(\br)\,\partial_n G_D(\br,\mathbf{y})\, dA,
  }
\end{equation}
where $f = \phi|_\Sigma$ is the prescribed Dirichlet data.  Both integrands are now
\emph{known}, so \eqref{eq:lec13:dirichlet-solution} is a genuine solution.

\textit{Significance:} The entire Dirichlet problem reduces to finding $G_D$.  Once $G_D$
is known for a given geometry, the solution for \emph{any} source $\rho$ and boundary
datum $f$ follows by a single integration.

\subsubsection*{Physical Interpretation of the $G_D$ Problem}
Equations \eqref{eq:lec13:GD-def} say: place a unit point charge at $\mathbf{y}$, remove
all other sources $\rho$, and hold the boundary $\Sigma$ at zero potential (grounded
conductor).  Solve for the resulting potential.  That potential is $G_D(\cdot,\mathbf{y})$.

This is precisely the \textbf{method of images} problem studied in earlier courses:
\begin{itemize}
  \item \textbf{Half-space:} a charge $q$ at height $d$ above a grounded infinite plane.
        Place an image charge $-q$ at $-d$; their combined potential vanishes on the plane
        and has the correct source singularity.  The result is $G_D$ for the half-space.
  \item \textbf{Sphere:} a charge $q$ at distance $r_0$ from the centre of a grounded
        conducting sphere of radius $a$.  The image charge is $-q\,a/r_0$ at distance
        $a^2/r_0$.  The resulting $G_D$ gives the solution for any source inside the sphere
        with any potential on the boundary.
\end{itemize}

\begin{remark}
  You used the method of images in year-1/2 courses without knowing you were computing
  Green's functions.  The new insight is that \emph{one} image calculation gives an
  \emph{infinite family} of solutions: plug different $\rho$ and $f$ into
  \eqref{eq:lec13:dirichlet-solution} and read off the answer.
\end{remark}

%------------------------------------------------------------
\subsection{The Neumann Green's Function}\label{subsec:lec13-neumann}

\subsubsection*{Motivation}
For the Neumann problem we know $\partial_n\phi|_\Sigma = g(\br)$ but not $\phi|_\Sigma$.
We want to eliminate the term $\phi\,\partial_n G$ from \eqref{eq:lec13:rep-general}.
The naive attempt is to demand $\partial_n G_N|_\Sigma = 0$.

\subsubsection*{Why the Naive Neumann Condition Fails}
Suppose $G_N$ satisfies \eqref{eq:lec13:G-def} and $\partial_n G_N|_\Sigma = 0$.
Integrate \eqref{eq:lec13:G-def} over the domain $\Omega$ (with $\mathbf{y}$ inside):
\begin{equation}\label{eq:lec13:naive-inconsistency}
  \int_\Omega \nabla^2_\br G_N\, d^3r = -4\pi
  \quad\xrightarrow{\text{Gauss}}\quad
  \oint_\Sigma \partial_n G_N\, dA = -4\pi.
\end{equation}
But if $\partial_n G_N|_\Sigma = 0$ then the left side is $0$.  This is a contradiction:
\begin{equation}\label{eq:lec13:contradiction}
  0 = -4\pi. \qquad \textit{(inconsistent!)}
\end{equation}
The naive condition is therefore \textbf{incompatible} with the defining equation.

\subsubsection*{The Correct Neumann Condition}
The flux integral must equal $-4\pi$.  The simplest consistent condition is to distribute
this equally over the boundary:
\begin{equation}\label{eq:lec13:GN-condition}
  \boxed{
    \partial_n G_N(\br,\mathbf{y})\big|_{\br\in\Sigma} = -\frac{4\pi}{S},
  }
\end{equation}
where $S = \oint_\Sigma dA$ is the total area of the boundary.

\dfn[dfn:green-neumann]{Neumann Green's Function $G_N$}{
  $G_N(\br,\mathbf{y})$ satisfies \eqref{eq:lec13:G-def} together with the boundary condition
  \eqref{eq:lec13:GN-condition}.
}

\subsubsection*{The Neumann Solution Formula}
Substituting $G_N$ into \eqref{eq:lec13:rep-general} and using \eqref{eq:lec13:GN-condition}:
\begin{equation}\label{eq:lec13:neumann-solution}
  \boxed{
    \phi(\mathbf{y})
    = \int_\Omega \rho(\br)\, G_N(\br,\mathbf{y})\, d^3r
    + \frac{1}{4\pi}\oint_\Sigma G_N(\br,\mathbf{y})\, g(\br)\, dA
    + \langle\phi\rangle_\Sigma,
  }
\end{equation}
where $g = \partial_n\phi|_\Sigma$ is the prescribed Neumann data and
$\langle\phi\rangle_\Sigma = \frac{1}{S}\oint_\Sigma\phi\,dA$ is the
(unknown) average value of $\phi$ on the boundary.

\textit{Significance:} The additive constant $\langle\phi\rangle_\Sigma$ is exactly the
ambiguity expected from uniqueness theory: the Neumann problem determines $\phi$ only up
to a global constant (physically, the potential reference level).  This constant can be
fixed by any convenient convention (e.g.\ setting $\langle\phi\rangle_\Sigma=0$).

%------------------------------------------------------------
\subsection{Comparison: Dirichlet vs.\ Neumann Problems}\label{subsec:lec13-comparison}

\begin{center}
\begin{tabular}{lll}
  \toprule
  & \textbf{Dirichlet} & \textbf{Neumann} \\
  \midrule
  Known data & $\phi|_\Sigma = f$ & $\partial_n\phi|_\Sigma = g$ \\
  Condition on $G$ & $G_D|_\Sigma = 0$ & $\partial_n G_N|_\Sigma = -4\pi/S$ \\
  Eliminated term & $G\,\partial_n\phi$ & $\phi\,\partial_n G$ \\
  Residual ambiguity & none & additive constant \\
  Image problem & unit charge at $\mathbf{y}$, grounded $\Sigma$ & unit charge at $\mathbf{y}$,
  $E_n = -4\pi/cS$ on $\Sigma$ \\
  \bottomrule
\end{tabular}
\end{center}

%------------------------------------------------------------
\subsection{Method of Images as Green's Function Computation}\label{subsec:lec13-images}

\subsubsection*{Conceptual Background}
Finding $G_D$ for a specific geometry is itself a boundary-value problem
(Poisson's equation with a point source, homogeneous Dirichlet condition), and it is in
general hard.  For highly symmetric geometries the method of images provides an elegant
exact answer.

\paragraph{Half-space ($z > 0$, grounded plane at $z=0$).}
Place a unit charge at $\mathbf{y} = (0,0,d)$.  The image is $-1$ at
$(0,0,-d)$.  Then
\begin{equation}\label{eq:lec13:GD-halfspace}
  G_D(\br,\mathbf{y}) = \frac{1}{|\br - \mathbf{y}|} - \frac{1}{|\br - \mathbf{y}^*|},
  \quad
  \mathbf{y}^* = (0,0,-d).
\end{equation}
One verifies $G_D|_{z=0} = 0$ (equal and opposite contributions cancel on the plane).

\paragraph{Sphere ($|\br| < a$, grounded at $|\br| = a$).}
For $\mathbf{y}$ at distance $r_0 < a$ from the centre, the image charge $-a/r_0$
is placed at $\mathbf{y}^* = (a^2/r_0)\hat{\mathbf{y}}$ (outside the sphere):
\begin{equation}\label{eq:lec13:GD-sphere}
  G_D(\br,\mathbf{y}) = \frac{1}{|\br - \mathbf{y}|} - \frac{a/r_0}{|\br - \mathbf{y}^*|}.
\end{equation}
One verifies $G_D|_{|\br|=a} = 0$.  This $G_D$ immediately gives the solution for
\emph{any} charge distribution inside the sphere with any Dirichlet boundary condition on
the sphere, via \eqref{eq:lec13:dirichlet-solution}.

\begin{figure}[h]
\centering
\begin{tikzpicture}[>=Stealth, thick, scale=1.1]
  % --- Half-space panel ---
  \begin{scope}[xshift=-4cm]
    % Grounded plane
    \draw[blue!60!black, thick] (-1.8,0) -- (1.8,0)
      node[right, blue!60!black] {$\Sigma:\; z=0$};
    \fill[blue!5] (-1.8,0) rectangle (1.8,-0.15);
    % Domain label
    \node at (0, 1.4) {$\Omega:\; z>0$};
    % Real charge
    \fill[red] (0, 1.0) circle (3pt) node[right] {$+1$ at $\mathbf{y}$};
    \draw[dashed] (0,1.0) -- (0,0) node[below] {\footnotesize $d$};
    % Image charge
    \fill[blue!70] (0, -1.0) circle (3pt) node[right] {$-1$ at $\mathbf{y}^*$};
    \draw[dashed] (0,-1.0) -- (0,0);
    \node at (0,-1.8) {\footnotesize Half-space: image at $-d$};
  \end{scope}
  % --- Sphere panel ---
  \begin{scope}[xshift=4cm]
    % Sphere boundary
    \draw[blue!60!black, thick] (0,0) circle (1.5);
    \node[blue!60!black] at (0,1.75) {$\Sigma:\;|\br|=a$};
    % Domain
    \node at (0,0.3) {$\Omega$};
    % Centre
    \fill (0,0) circle (1.5pt);
    % Real charge y inside
    \fill[red] (0.8, 0) circle (3pt) node[above] {$+1$};
    \node[red, below] at (0.8,0) {$\mathbf{y}$};
    \draw[|<->|, gray] (0,0.05) -- (0.8,0.05) node[midway,above]{\footnotesize $r_0$};
    % Image charge outside
    \fill[blue!70] (2.8, 0) circle (3pt) node[right] {$-\frac{a}{r_0}$};
    \node[below, blue!70] at (2.8,0) {$\mathbf{y}^*$};
    \draw[|<->|, gray] (0,-0.1) -- (2.8,-0.1) node[midway,below]{\footnotesize $a^2/r_0$};
    \node at (0,-2.0) {\footnotesize Sphere: image at $a^2/r_0$};
  \end{scope}
\end{tikzpicture}
\caption{Method of images as computation of the Dirichlet Green's function.
\textit{Left}: grounded half-space; image charge $-1$ at $\mathbf{y}^*=(0,0,-d)$.
\textit{Right}: grounded sphere of radius $a$; image charge $-a/r_0$ at distance $a^2/r_0$.
In both cases the image ensures $G_D=0$ on $\Sigma$.}
\label{fig:lec13:images}
\end{figure}

%------------------------------------------------------------
\subsection{Symmetry of the Dirichlet Green's Function}\label{subsec:lec13-symmetry}

\subsubsection*{Mathematical Setup}
It is not obvious from the definition \eqref{eq:lec13:GD-def} that
$G_D(\br,\mathbf{y}) = G_D(\mathbf{y},\br)$.
The free-space piece $1/|\br-\mathbf{y}|$ is manifestly symmetric.  The harmonic part
$F$ that enforces the boundary condition is not obviously so.

For the half-space and sphere examples above one can check symmetry by direct calculation.
In general:

\clm[clm:GD-symmetry]{Symmetry of $G_D$}{
  The Dirichlet Green's function satisfies
  \begin{equation}\label{eq:lec13:GD-sym}
    \boxed{G_D(\br, \mathbf{y}) = G_D(\mathbf{y}, \br)}
  \end{equation}
  for any $\br, \mathbf{y} \in \Omega$.
}

\textit{Proof (sketch, to be completed next lecture):} Apply Green's second identity with
$u = G_D(\cdot,\mathbf{y})$ and $v = G_D(\cdot,\mathbf{z})$.  The boundary terms vanish
(since $G_D = 0$ on $\Sigma$); the volume integrals yield
$G_D(\mathbf{z},\mathbf{y}) - G_D(\mathbf{y},\mathbf{z}) = 0$.

\textit{Physical meaning:} The potential at $\mathbf{y}$ due to a unit charge at $\br$
(with grounded boundary) equals the potential at $\br$ due to a unit charge at $\mathbf{y}$.
This is a reciprocity theorem for electrostatics.

%------------------------------------------------------------
\subsection{Summary}\label{subsec:lec13-summary}

\begin{itemize}
  \item \textbf{Mean value property:} $u(\mathbf{y}) = \frac{1}{4\pi R^2}\oint_{|\br-\mathbf{y}|=R} u\, dA$
        for any harmonic $u$ and any $R$ with $B_R(\mathbf{y})\subset\Omega$.
  \item \textbf{Earnshaw:} harmonic functions have no interior extrema (follows immediately
        from mean value property).
  \item \textbf{General Green's function:} $G = 1/|\br-\mathbf{y}| + F$ where $F$ is
        harmonic; $G$ satisfies $\nabla^2_\br G = -4\pi\delta(\br-\mathbf{y})$.
  \item \textbf{Dirichlet:} $G_D|_\Sigma = 0$; gives
        $\phi(\mathbf{y}) = \int_\Omega\rho G_D\,d^3r
        - \frac{1}{4\pi}\oint_\Sigma f\,\partial_n G_D\,dA$.
        Physical interpretation: $G_D$ is the potential of a unit charge at $\mathbf{y}$
        in a grounded domain — precisely the method-of-images problem.
  \item \textbf{Neumann:} $\partial_n G_N|_\Sigma = -4\pi/S$ (not zero — that is
        inconsistent); gives $\phi$ up to an additive constant $\langle\phi\rangle_\Sigma$.
  \item \textbf{Symmetry:} $G_D(\br,\mathbf{y}) = G_D(\mathbf{y},\br)$ (proof next lecture).
\end{itemize}
